\documentclass{article}
\usepackage[utf8]{inputenc}
\usepackage[russian]{babel}
\usepackage{amsmath}
\usepackage{graphicx}
\usepackage{float}
\usepackage{caption}
\usepackage{subcaption}
\usepackage{geometry}

\geometry{
 a4paper,
 total={170mm,257mm},
 left=20mm,
 top=20mm,
}

\title{Исследование динамики стохастической системы второго порядка:\\ эффект усиленной шумом стабильности (NES)}
\date{\today}

\begin{document}

\maketitle

\section{Введение и теоретическая модель}
Динамика многих физических, биологических и радиотехнических систем (например, джозефсоновских контактов, кольцевых лазеров, маятников с вязким трением) с учетом диффузионных флуктуаций может быть описана стохастическим дифференциальным уравнением (СДУ) Ланжевена второго порядка:
\begin{equation}
    \ddot{x} + \alpha \dot{x} + \sin(x) = a + A \sin(\omega t) + \xi(t),
    \label{eq:main}
\end{equation}
где:
\begin{itemize}
    \item $x(t)$ — координата системы (фазовая переменная);
    \item $\alpha$ — коэффициент затухания (вязкого трения), определяющий диссипацию энергии;
    \item $a$ — постоянная составляющая внешней силы, задающая наклон профиля потенциала;
    \item $A, \omega$ — амплитуда и частота гармонического переменного сигнала;
    \item $\xi(t)$ — белый гауссов шум с нулевым математическим ожиданием $\langle \xi(t) \rangle = 0$ и корреляционной функцией $\langle \xi(t)\xi(t+\tau) \rangle = D \delta(\tau)$, где $D$ — интенсивность шума.
\end{itemize}

Эффективный потенциал системы имеет вид:
\begin{equation}
    U(x) = -a x - \cos(x).
\end{equation}
Цель работы — проанализировать среднее время переключения (Mean Switching Time, MST), соответствующее переходу частицы из одной потенциальной ямы в соседнюю (достижение границы $x_{th} = \pi \approx 3.14$) в зависимости от интенсивности шума $D$ при различных режимах затухания $\alpha$.

\section{Численный метод Хюна}
Для высокой точности интегрирования осциллирующей компоненты СДУ (\ref{eq:main}) разбивается на систему уравнений и решается стохастическим методом Хюна. Это двухэтапный метод предиктора-корректора, который обеспечивает лучшую сходимость по сравнению с методом Эйлера. Среднее время MST вычисляется параллельно из 1000 реализаций.


\newpage
\section{Результаты моделирования}
Для анализа эффекта усиленной шумом стабильности параметры системы были зафиксированы в конфигурации: $a = 0.5$, $A = 1.0$, $\omega = 0.4$. В исследовании рассматриваются три режима трения $\alpha$. Интенсивность шума $D$ варьируется логарифмически от $10^{-3}$ до $10^0$.

\subsection{Промежуточный режим ($\alpha = 1.0$): Эффект NES}
В этом режиме установлен оптимальный баланс между накачкой энергии и диссипацией. 

\begin{figure}[H]
    \centering
    \includegraphics[width=0.85\textwidth]{alpha_1.png} % МЕСТО ДЛЯ ГРАФИКА (Скриншот 1)
    \caption{Оптимальное затухание $\alpha = 1.0$. Наблюдается эффект Noise Enhanced Stability.}
    \label{fig:alpha_1}
\end{figure}

При исчезающе малом шуме детерминированный сигнал синхронно выталкивает частицу за барьер за минимальное время (MST $\approx 6.0$ с). Однако при увеличении шума (в диапазоне от $10^{-2}$ до $10^{-1}$) наблюдается аномальный рост времени задержки. Шум играет роль десинхронизатора: он сбивает фазу частицы во время оптимального «окна» для перескока, заставляя систему оставаться в потенциальной яме еще на несколько периодов сигнала.

\newpage
\subsection{Инерционный режим слабого трения ($\alpha = 0.2$)}
Среда обладает малой вязкостью, система накапливает кинетическую энергию (underdamped).

\begin{figure}[H]
    \centering
    \includegraphics[width=0.85\textwidth]{alpha_02.png} % МЕСТО ДЛЯ ГРАФИКА (Скриншот 3)
    \caption{Режим слабого трения $\alpha = 0.2$. Инерционный пролет барьеров.}
    \label{fig:alpha_02}
\end{figure}

Барьеры практически не останавливают частицу — детерминированная траектория пересекает предельное значение крайне быстро (MST $\approx 3.8$ секунд). Влияние шума здесь сводится лишь к небольшой десинхронизации, которая статистически слегка замедляет систему, но качественных пиков задержки не образуется из-за огромного инерционного импульса.

\newpage
\subsection{Режим сильного (вязкого) трения ($\alpha = 5.0$)}
Динамика соответствует передемпфированному режиму. Вязкость среды поглощает энергию быстрее, чем внешний сигнал успевает ее накачать.

\begin{figure}[H]
    \centering
    \includegraphics[width=0.85\textwidth]{alpha_5.png} % МЕСТО ДЛЯ ГРАФИКА (Скриншот 2)
    \caption{Режим сильного трения $\alpha = 5.0$. Система заперта.}
    \label{fig:alpha_5}
\end{figure}

Суммарного воздействия $a + A$ недостаточно для преодоления барьера из-за потерь. Частица совершает вынужденные осцилляции на дне ямы, никогда не достигая порога за отведенное время симуляции ($t=1000$). Даже при максимальном в рамках эксперимента шуме $D = 1.0$, вероятности Крамерсовского (термоактивационного) перескока недостаточно, что отражается в плоском графике на уровне максимального времени наблюдения (MST = 1000).

\section{Заключение}
В ходе экспериментов было продемонстрировано неожиданное явление нелинейной стохастической динамики — эффект усиленной шумом стабильности (NES). Показано, что в промежуточном режиме трения ($\alpha=1.0$) добавление шума не разрушает, а напротив — стабилизирует метастабильное состояние частицы в потенциальной яме, увеличивая среднее время её удержания почти вдвое за счет десинхронизации детерминированного периодического механизма переброса.

\end{document}
